\chapter{Inexact Newton and forcing terms}
\label{ch:inexact}
\impl{src/linear\_solvers.jl}

\section{The over-solving problem}

Newton's step \eqref{eq:newton} is a linear system, and when it is solved by an
iterative method the accuracy of that solve is a free parameter.  An
\emph{inexact Newton method} accepts any $\Delta\Uvec_k$ satisfying
\begin{equation}
  \norm{\Rvec(\Uvec_k) - \Kmat(\Uvec_k)\Delta\Uvec_k} \le \eta_k \norm{\Rvec(\Uvec_k)},
  \label{eq:inexact}
\end{equation}
where $\eta_k \in [0,1)$ is the \emph{forcing term}.  Choosing $\eta_k$
constant and small --- the naive choice, and Carina's original behaviour ---
solves every linear system to the same tolerance regardless of how far the
nonlinear iteration still has to go.

That this wastes work is not an abstract concern.  On a $530{,}523$-DOF Newmark
step, conjugate gradients spent $181$, $205$, and $193$ iterations on the three
solves of a step while the nonlinear residual fell
\begin{equation*}
  1.61\times10^{4} \ \longrightarrow\ 1.34\times10^{2}
  \ \longrightarrow\ 7.79\times10^{-2} \ \longrightarrow\ 4.60\times10^{-8}.
\end{equation*}
The first two steps were computed to eight digits of linear accuracy and then
discarded by the next Newton correction.  The pattern on a quasi-static $J_2$
specimen is worse: roughly $350$ iterations on each of the first seven solves
while the residual crawls from $1.57\times10^3$ to $1.76\times10^2$, after
which the line search backtracks to $\alpha = \tfrac14$ and throws away most of
the step it just paid for.

\section{Eisenstat--Walker choice 2}

The remedy is to tie $\eta_k$ to the convergence actually being observed.
Carina implements choice~2 of Eisenstat and Walker~\cite{EisenstatWalker1996}:
\begin{equation}
  \eta_k^{\mathrm{EW}} = \gamma \left(\frac{\norm{\Rvec_k}}{\norm{\Rvec_{k-1}}}\right)^{\alpha},
  \qquad \gamma \in (0,1], \quad \alpha \in (1,2].
  \label{eq:ew}
\end{equation}
The reasoning is that the ratio of successive residual norms estimates how well
the linear model predicted the nonlinear decrease.  When the previous step
behaved as the model said it would, the ratio is small and a tighter solve is
warranted; when the model is doing badly, a tight solve is refining a direction
that is wrong anyway.  The default $\alpha = (1+\sqrt5)/2$ is the paper's, and
recovers superlinear convergence of the outer iteration.

On the first iteration of a step no ratio exists, and $\eta_0$ is taken from
the \code{initial} parameter.

\section{Three safeguards}
\label{sec:ew-safeguards}

Equation~\eqref{eq:ew} alone is not usable.  Carina composes it with three
guards, applied in order.

\paragraph{1. Against $\eta$ collapsing too fast.}  A single unusually good
step can drive \eqref{eq:ew} to a very small value that the local convergence
rate does not justify, wasting on the next solve exactly what the method exists
to save.  Eisenstat and Walker's own safeguard raises $\eta$ back toward its
predecessor whenever that predecessor was itself large:
\begin{equation}
  \eta_k \leftarrow \max\left(\eta_k^{\mathrm{EW}},\ \gamma\,\eta_{k-1}^{\alpha}\right)
  \qquad\text{whenever}\qquad \gamma\,\eta_{k-1}^{\alpha} > 0.1 .
  \label{eq:ew-guard}
\end{equation}
The threshold $0.1$ is the paper's and is not a tuning parameter in Carina.  It
nonetheless has a consequence for the parameter that \emph{is} tunable, which
Section~\ref{sec:etamax} takes up.

\paragraph{2. Against over-solving the final step.}  Near convergence the ratio
in \eqref{eq:ew} becomes tiny and demands far more linear accuracy than the
nonlinear stopping test will ever reward.  Kelley's guard~\cite{Kelley1995}
bounds $\eta$ from below by what the target actually requires:
\begin{equation}
  \eta_k \leftarrow \max\left(\eta_k,\ \frac{s\,\tau}{\norm{\Rvec_k}}\right),
  \label{eq:kelley}
\end{equation}
where $\tau$ is the residual norm the nonlinear test is aiming at and $s$ a
safety factor, $\tfrac12$ by default.  Carina reads $\tau$ from the status test
tree of Section~\ref{sec:status}, so a deck's own \code{termination} block
governs rather than a nominal tolerance it may have overridden.

\paragraph{3. A hard clamp.}  Finally
\begin{equation}
  \eta_k \leftarrow \min\left(\max\left(\eta_k,\ \varepsilon_{\mathrm{lin}}\right),\ \eta_{\max}\right),
  \label{eq:clamp}
\end{equation}
where $\varepsilon_{\mathrm{lin}}$ is the deck's own linear tolerance.

The order matters, and Algorithm~\ref{alg:forcing} fixes it: the two lower
bounds are applied before the clamp, so neither can push $\eta$ outside
$[\varepsilon_{\mathrm{lin}}, \eta_{\max}]$.

\begin{algorithm}[h]
\caption{Forcing term at Newton iteration $k$}
\label{alg:forcing}
\begin{algorithmic}[1]
  \Require $\norm{\Rvec_k}$; previous $\norm{\Rvec_{k-1}}$ and $\eta_{k-1}$
           if any; target $\tau$; parameters $\gamma, \alpha, \eta_{\max},
           \eta_0, s$; linear tolerance $\varepsilon_{\mathrm{lin}}$
  \If{no previous residual}
    \State $\eta \gets \eta_0$
      \Comment{first iteration of a step}
  \Else
    \State $\eta \gets \gamma\left(\norm{\Rvec_k}/\norm{\Rvec_{k-1}}\right)^{\alpha}$
    \If{$\gamma\,\eta_{k-1}^{\alpha} > 0.1$}
      \State $\eta \gets \max\left(\eta,\ \gamma\,\eta_{k-1}^{\alpha}\right)$
        \Comment{safeguard 1: anti-collapse}
    \EndIf
  \EndIf
  \If{$s > 0$ \textbf{and} $\norm{\Rvec_k} > 0$}
    \State $\eta \gets \max\left(\eta,\ s\,\tau/\norm{\Rvec_k}\right)$
      \Comment{safeguard 2: over-solve guard}
  \EndIf
  \State $\eta_k \gets \min\left(\max\left(\eta,\ \varepsilon_{\mathrm{lin}}\right),\ \eta_{\max}\right)$
    \Comment{safeguard 3: clamp}
  \State \Return $\eta_k$
\end{algorithmic}
\end{algorithm}

\begin{proposition}[Safety]
\label{prop:safety}
Under \eqref{eq:clamp}, a forcing term can only loosen a linear solve relative
to the fixed-tolerance baseline, never tighten one.
\end{proposition}

\noindent This is the property that makes the method switchable on without
re-validating a model.  The lower clamp guarantees the final Newton iterations
run at exactly the tolerance the user requested, so the converged solution is
as accurate as it was before and an A/B comparison against a fixed tolerance
compares like with like.  Measurements bear this out: across CPU and A100 runs
in both implicit regimes, $\norm{\Uvec}_\infty$ agrees with the
fixed-tolerance baseline to every printed digit at every output stop.

\begin{remark}
Preconditioned conjugate gradients stops on the preconditioned residual
$\Prec$-norm relative to its own initial value, not on the unpreconditioned
ratio of \eqref{eq:inexact} for which the theory is stated.  The forcing term
is a heuristic in either norm; Proposition~\ref{prop:safety} is what bounds the
damage when the two disagree.
\end{remark}

\section{Choosing \texorpdfstring{$\eta_{\max}$}{eta max}}
\label{sec:etamax}

The literature's default is $\eta_{\max} = 0.9$.  Carina's is $0.2$, and the
reason is structural rather than empirical fitting.

Measurements across a CPU quasi-static sweep and two A100 regimes show that the
\emph{total} linear work is nearly invariant in $\eta_{\max}$: $63.1$k, $62.9$k
and $63.6$k CG iterations at $\eta_{\max} = 0.1$, $0.2$, $0.5$ respectively on
the $J_2$ specimen, and $1019$ versus $1011$ on the A100 Newmark problem.  What
$\eta_{\max}$ actually controls is how many \emph{Newton} iterations are spent
buying that work --- $413$ versus $547$ on the same CPU sweep.

The mechanism is safeguard~\eqref{eq:ew-guard}, which activates when
$\gamma\eta^{\alpha} > 0.1$.  At the default $\alpha$:
\begin{center}
\begin{tabular}{@{}lll@{}}
\toprule
$\eta_{\max}$ & $\eta_{\max}^{\alpha}$ & Behaviour \\
\midrule
$0.2$ & $0.076$ & below the knee; the residual ratio drives $\eta$ \\
$0.5$ & $0.326$ & above it; the safeguard pins $\eta$ high for several iterations \\
$0.9$ & $0.842$ & far above \\
\bottomrule
\end{tabular}
\end{center}
So $0.2$ is the largest round value that keeps Eisenstat--Walker's
anti-collapse safeguard out of the way of Eisenstat--Walker's own adaptivity.
Two different problems on two different architectures selected it
independently.

\section{What the method buys}

A forcing term trades linear work for nonlinear work.  It therefore pays in
proportion to how much of a step is the linear solve --- most on the
matrix-free GPU paths of Chapter~\ref{ch:matrixfree}, where the step
essentially \emph{is} the stiffness matrix--vector product, and least on
assembled CPU runs where residual assembly and the line search already
dominate.  Measured with the defaults:

\begin{center}
\begin{tabular}{@{}lcc@{}}
\toprule
Configuration & CG iterations & Per-step wall \\
\midrule
CPU quasi-static, $J_2$, $57$k DOF        & $2.47\times$ fewer & $1.30\times$ \\
A100 Newmark, $530$k DOF, CG + Jacobi     & $2.27\times$ fewer & $1.45\times$ \\
A100 quasi-static, $530$k DOF, CG + AMG   & $2.05\times$ fewer & $1.59\times$ \\
\bottomrule
\end{tabular}
\end{center}

The quasi-static figure is on top of the roughly twofold advantage algebraic
multigrid already holds over Jacobi on that problem.

\section{Interaction with lagged preconditioners}

The hierarchy staleness detector of Section~\ref{sec:amg-lagging} compares
conjugate-gradient iteration counts across solves.  Under a forcing term those
counts are taken at \emph{different} tolerances and are not comparable as raw
numbers.

Simply ignoring loosened solves does not work, and the reason is
guard~\eqref{eq:kelley}: it can raise $\eta$ enough that a run converges
without ever solving at the deck's tolerance, so a detector that only counted
full-tolerance solves would never fire at all.  Since that detector is the only
rebuild trigger on the quasi-static path, the hierarchy would be built once at
the reference configuration and kept for an entire load history --- precisely
the failure Section~\ref{sec:nullspace} exists to prevent.

Carina instead rescales.  Because a Krylov residual falls geometrically, the
iteration count is roughly proportional to the number of digits of reduction
requested, so
\begin{equation}
  \hat{m}_k = m_k \cdot \frac{\log \varepsilon_{\mathrm{lin}}}{\log \eta_k}
  \label{eq:rescale}
\end{equation}
converts a count $m_k$ taken at tolerance $\eta_k$ into the count the same
solve would have cost at the deck's tolerance.  At $\eta_k =
\varepsilon_{\mathrm{lin}}$ this is exactly the identity, so a run with no
forcing term feeds the detector precisely the numbers it always did.  The
estimate is coarse, but the detector's threshold --- a factor of three ---
is coarser.
